\documentclass{article}
\usepackage[UTF8]{ctex}
\usepackage{amsmath,amssymb,amsfonts}
\usepackage{enumitem}
\usepackage{geometry}
\geometry{a4paper, margin=1in}
% === 错题本排版优化（不参与题目标识） ===
\setlength{\parindent}{0pt}
\setlength{\parskip}{0.35em plus 0.10em minus 0.08em}
\setlength{\emergencystretch}{3em}
\setlist{itemsep=0.25em, topsep=0.40em, parsep=0pt}
\allowdisplaybreaks[2]
% === 错题本排版优化结束 ===
\begin{document}
\section*{湘潭大学2024年数学分析真题}
\section*{196．湘潭大学2024年数学分析试卷}

\subsection{题目1}
一．计算题．每题10分，共30分．

\begin{itemize}
  \item[（1）] 求极限 \(\lim _{x \rightarrow 0} \frac{\tan x-\sin x}{x^{3}}\) ．
  \item[（2）] 求极限 \(\lim _{n \rightarrow \infty} \int_{0}^{1} \frac{x^{n}}{1+x} \mathrm{~d} x\) ．
  \item[（3）] 求极限 \(\lim _{x \rightarrow+\infty} x^{\frac{1}{x}}\) ．
\end{itemize}

\subsection{题目2}
二．计算题．每题10分，共20分．

\begin{itemize}
  \item[（1）] 设 \(y=\frac{\sqrt{1+x^{2}}}{x}\) ，求 \(\mathrm{d}^{2} y\) 。
  \item[（2）] 设 \(I(y)=\int_{y}^{y^{2}} \frac{\cos x y}{x} \mathrm{~d} x\) ，求 \(I^{\prime}(y)\) ．
\end{itemize}

\subsection{题目3}
三。（15 分）设 \(\Omega\) 为球面 \(x^{2}+y^{2}+z^{2}=3\) 和抛物面 \(2 z=x^{2}+y^{2}\) 所围成的区域，计算重积分

\[
\iiint_{\Omega}(x+y+z)^{2} \mathrm{~d} x \mathrm{~d} y \mathrm{~d} z
\]

\subsection{题目4}
四．（15 分）已知 \(\sum\) 为封闭曲面，l 为固定向量， \(\mathbf{n}\) 为 \(\sum\) 的外法向量，证明： \(\iint_{\sum} \cos (\mathbf{n}, \mathbf{1}) \mathrm{d} S=0\) ．

\subsection{题目5}
五．（20 分）已知含参量积分 \(I=\int_{0}^{+\infty} e^{-\alpha x} \sin x \mathrm{~d} x\) ．

\begin{itemize}
  \item[（1）] 对于 \(\alpha \geq \alpha_{0}>0\) ，讨论 \(I\) 的一致收敛性．
  \item[（2）] 对于 \(\alpha>0\) ，讨论 \(I\) 的一致收敛性．
\end{itemize}

\subsection{题目6}
六．（15 分）叙述 \(\lim _{x \rightarrow x_{0}^{+}} f(x)\) 存在且有限的柯西收敛准则并证明．

\subsection{题目7}
七．（15 分）将 \(f(x)=x^{2}\) 在 \([0,2 \pi]\) 上傅里叶展开．

\subsection{题目8}
八．（20 分）设 \(\zeta(s)=\sum_{n=1}^{\infty} \frac{1}{n^{s}}, E(s)=\prod_{n=1}^{\infty} \frac{1}{1-p_{n}^{s}}, s>1\) ，其中 \(p_{n}\) 为第 \(n\) 个质数。

\begin{itemize}
  \item[（1）] 证明：\(\zeta(s)\) 与 \(E(s)\) 均收敛。
  \item[（2）] 证明：\(\zeta(s)=E(s)\) ．
\end{itemize}
\end{document}
